\section{Related Work}
\label{sec:related}

\para{LLM scientific ideation and its weaknesses.} The pivotal empirical result is \citet{si2024can}, who ran a blind study with 49 expert idea writers and 79 reviewers across 7 NLP topics. Their LLM agent---RAG over retrieved papers, over-generation of thousands of seeds, deduplication, and pairwise tournament ranking---produced ideas judged \emph{more novel} than expert humans, but slightly less feasible. The study also exposed two systemic weaknesses that frame much of the follow-up work: \emph{diversity collapse} (thousands of generated seeds yield only a few hundred unique ideas) and \emph{unreliable self-evaluation} (the LLM ranker barely exceeded chance consistency). We reuse their 7 topics and 1--10 rubric so our numbers are comparable, and we directly test whether grounding mitigates diversity collapse.

\para{Grounding via literature comparison.} SciMON~\citep{wang2023scimon} retrieves ``inspirations,'' generates, then iteratively compares each candidate to prior work and revises until it is novel enough. Chain-of-Ideas~\citep{li2024chain} organizes literature into a progression chain and ships a pairwise ``Idea Arena'' evaluation, and Nova~\citep{nova2024} uses planned iterative retrieval to produce more unique top-rated ideas. These systems ground novelty against \emph{retrieved prior work}. Our \litcomp condition implements exactly this retrieve-and-differentiate loop, and our objective \lgn metric is the same idea applied to \emph{evaluation}: distance to retrieved prior work rather than an LLM's opinion.

\para{Grounding via induction and deduction.} MOOSE~\citep{yang2023moose} performs hypothetical induction from a raw corpus to novel social-science hypotheses, refined by feedback. \citet{causalgraph2024} build a causal knowledge graph from tens of thousands of psychology articles and show that an LLM grounded in this inductive structure beats an LLM-only baseline on novelty. On the deductive side, HypoGen~\citep{hypogen2025} introduces the Bit$\rightarrow$Flip$\rightarrow$Spark schema---state the prevailing assumption, invert it, derive the idea---which we adopt verbatim as our \deduction scaffold; HypoSpace~\citep{hypospace2025} evaluates LLMs as set-valued hypothesis generators. Unlike these works, which each validate a single mechanism on a bespoke benchmark, we place induction and deduction on the same footing as four other mechanisms.

\para{Grounding via proposals, outcomes, and experiments.} \citet{dynamiccontrol2024} train SFT+RL reward models for novelty, feasibility, and effectiveness, defining the rubric conventions we follow. VirSci~\citep{virsci2024} uses multi-agent teams to generate, evaluate, and refine full proposals. AutoDiscovery~\citep{autodiscovery2025} drives open-ended discovery by Bayesian surprise---the epistemic shift after running experiments---which motivates our \outcome scaffold. Deep-Ideation~\citep{deepideation2025} evolves ideas over a concept network using a critic trained on real reviewer feedback, which we operationalize as our \smallexp scaffold (design a falsification test, predict its result, self-critique, revise). Because generic ideas cannot be cheaply executed in-loop, we follow the predict-and-critique operationalization used by MOOSE and Deep-Ideation.

\para{Grounding the evaluation.} Two recent papers warn that the metric, not just the generator, must be grounded. \emph{All That Glitters Is Not Novel}~\citep{glitters2025} finds that a quarter of ``novel'' AI research documents are smartly borrowed and bypass detectors, recommending that novelty be measured as distance to retrieved prior work. \emph{On the Limits of LLM-as-Judge}~\citep{rqbench2026} documents a \emph{novelty mirage}: LLM judges rate model-generated questions as highly novel while human experts prefer the real references. These cautions motivate our three-lens evaluation and the convergent-validity test that is, to our knowledge, the first controlled measurement of the mirage on a single fixed idea set. Surveys of the area~\citep{survey2025a,survey2025b} report a directional ranking (graph-grounded $3.83 >$ Chain-of-Ideas $3.21 >$ vanilla $2.44$) suggesting grounding helps novelty; our objective metric and overall-quality scores caution that this ranking may itself be judge-inflated.
